% Options for packages loaded elsewhere
\PassOptionsToPackage{unicode}{hyperref}
\PassOptionsToPackage{hyphens}{url}
\PassOptionsToPackage{dvipsnames,svgnames,x11names}{xcolor}
%
\documentclass[
  10pt,
]{article}
\usepackage{amsmath,amssymb}
\usepackage{iftex}
\ifPDFTeX
  \usepackage[T1]{fontenc}
  \usepackage[utf8]{inputenc}
  \usepackage{textcomp} % provide euro and other symbols
\else % if luatex or xetex
  \usepackage{unicode-math} % this also loads fontspec
  \defaultfontfeatures{Scale=MatchLowercase}
  \defaultfontfeatures[\rmfamily]{Ligatures=TeX,Scale=1}
\fi
\usepackage{lmodern}
\ifPDFTeX\else
  % xetex/luatex font selection
  \setmainfont[]{DejaVu Serif}
  \setsansfont[]{DejaVu Sans}
  \setmonofont[]{DejaVu Sans Mono}
  \setmathfont[]{Latin Modern Math}
\fi
% Use upquote if available, for straight quotes in verbatim environments
\IfFileExists{upquote.sty}{\usepackage{upquote}}{}
\IfFileExists{microtype.sty}{% use microtype if available
  \usepackage[]{microtype}
  \UseMicrotypeSet[protrusion]{basicmath} % disable protrusion for tt fonts
}{}
\makeatletter
\@ifundefined{KOMAClassName}{% if non-KOMA class
  \IfFileExists{parskip.sty}{%
    \usepackage{parskip}
  }{% else
    \setlength{\parindent}{0pt}
    \setlength{\parskip}{6pt plus 2pt minus 1pt}}
}{% if KOMA class
  \KOMAoptions{parskip=half}}
\makeatother
\usepackage{xcolor}
\usepackage[margin=2.0cm]{geometry}
\usepackage{longtable,booktabs,array}
\usepackage{calc} % for calculating minipage widths
% Correct order of tables after \paragraph or \subparagraph
\usepackage{etoolbox}
\makeatletter
\patchcmd\longtable{\par}{\if@noskipsec\mbox{}\fi\par}{}{}
\makeatother
% Allow footnotes in longtable head/foot
\IfFileExists{footnotehyper.sty}{\usepackage{footnotehyper}}{\usepackage{footnote}}
\makesavenoteenv{longtable}
\setlength{\emergencystretch}{3em} % prevent overfull lines
\providecommand{\tightlist}{%
  \setlength{\itemsep}{0pt}\setlength{\parskip}{0pt}}
\setcounter{secnumdepth}{-\maxdimen} % remove section numbering
% doc_header.tex -- LaTeX preamble additions used by build_docs.sh for every
% Markdown -> LaTeX -> PDF conversion (the documents of the folder Revision_Lord_Gravity).
\usepackage{etoolbox}
% longtable may not be loaded by pandoc when a document has no table
\usepackage{longtable}
% bold vectors: \bm -> \symbf under unicode-math (XeLaTeX), bm package otherwise
\makeatletter
\@ifpackageloaded{unicode-math}{\providecommand{\bm}[1]{\symbf{#1}}}{\usepackage{bm}}
\makeatother
% smaller type in long tables so that wide audit tables fit the text width;
% inside a table a line break is allowed after an underscore of plain text
% (pandoc writes it as \_), so that long identifiers wrap inside their column
% rows are set a little apart, so that symbols with indices do not touch
\AtBeginEnvironment{longtable}{\footnotesize\renewcommand{\arraystretch}{1.45}%
  \let\sptunderscore\_\def\_{\sptunderscore\allowbreak}}
\setlength{\LTcapwidth}{\textwidth}
% A short table is kept on one page.  The build (code/scripts/keep_tables.py) ends the rows of such a table with \\*,
% which forbids a page break after a row, and sets \sptkeptpenalty to 10000 in front of the table.  This penalty is
% put here in the place of the breakpoint (penalty 0) that longtable leaves after every table, so that the paragraph
% that follows a kept table, its caption, begins on the same page; it is set back to 0 after the table.
\newcount\sptkeptpenalty
\sptkeptpenalty=0
\makeatletter
\patchcmd{\endlongtable}{\endgraf\penalty\z@\addvspace\LTpost}%
  {\endgraf\penalty\sptkeptpenalty\addvspace\LTpost\global\sptkeptpenalty=\z@}%
  {}{\PackageError{doc_header}{the end of longtable could not be patched}{The tables can still be typeset, but a short table may be separated from its caption.}}
\makeatother
% allow line breaks after slashes and hyphens inside table cells
\usepackage{xurl}
% code blocks: break over-long lines instead of running into the margin.
% Plain ``` blocks are typeset by pandoc as {verbatim}, highlighted ones
% (```bash) as {Highlighting}; both are re-declared as breakable fancyvrb
% environments.  The continuation line is marked with a small hooked arrow.
\usepackage{fvextra}
\RecustomVerbatimEnvironment{verbatim}{Verbatim}{breaklines,breakanywhere,fontsize=\small}
\makeatletter
\@ifundefined{Highlighting}{}{%
  \RecustomVerbatimEnvironment{Highlighting}{Verbatim}{breaklines,breakanywhere,fontsize=\small,commandchars=\\\{\}}}
\makeatother
% inline code (pandoc writes it as \texttt{...}): never hyphenate -- an inserted
% hyphen could be read as part of a file name -- but allow a line break after
% an underscore, slash, comma, colon or equals sign, so that long paths and
% column lists wrap inside the text width and inside table columns.
% Inside a formula (e.g. the name of a notebook symbol in an equation) the same
% command must first leave math mode, which \text of amsmath does.
\ExplSyntaxOn
\tl_new:N \l__spt_code_tl
\cs_new_protected:Npn \__spt_texttt:n #1
  {
    \tl_set:Nn \l__spt_code_tl {#1}
    \regex_replace_all:nnN { ([\/\,\:\=]) } { \1 \c{allowbreak} } \l__spt_code_tl
    \regex_replace_all:nnN { \c{_} } { \c{_} \c{allowbreak} } \l__spt_code_tl
    \group_begin: \ttfamily \hyphenchar \font = -1 \scan_stop: \l__spt_code_tl \group_end:
  }
\RenewDocumentCommand \texttt { m }
  {
    \mode_if_math:TF { \text { \__spt_texttt:n {#1} } } { \__spt_texttt:n {#1} }
  }
\ExplSyntaxOff
% A figure is written in the Markdown source as an image paragraph followed by
% a caption paragraph ("**Figure S1.** ..."), so that the caption is also
% visible on GitHub.  An image that starts a paragraph is centered.
\usepackage{graphicx}% pandoc loads it only for documents that contain an image
\NewCommandCopy\sptincludegraphics\includegraphics
\RenewDocumentCommand\includegraphics{O{} m}{%
  \ifvmode\begin{center}\sptincludegraphics[#1]{#2}\end{center}%
  \else\sptincludegraphics[#1]{#2}\fi}
% keep figures where they are referenced
\usepackage{float}
\floatplacement{figure}{H}
% running header with the document title
\usepackage{fancyhdr}
\pagestyle{fancy}
\fancyhf{}
\fancyhead[L]{\small PROJECT spacetimes, Revision\_Lord\_Gravity}
\fancyhead[R]{\small\thepage}
\renewcommand{\headrulewidth}{0.4pt}
% 16x16 matrices: tighter columns so that a full matrix fits the text width
\setlength{\arraycolsep}{3.2pt}
% The documents of this project contain long tagged symbols, file paths and
% identifiers that cannot be hyphenated.  A paragraph is therefore allowed to
% have looser lines before TeX gives up and lets a line stick out of the margin.
\AtBeginDocument{\tolerance=2000 \emergencystretch=5em \hbadness=2000}
\ifLuaTeX
  \usepackage{selnolig}  % disable illegal ligatures
\fi
\IfFileExists{bookmark.sty}{\usepackage{bookmark}}{\usepackage{hyperref}}
\IfFileExists{xurl.sty}{\usepackage{xurl}}{} % add URL line breaks if available
\urlstyle{same}
\hypersetup{
  pdftitle={The remarks of the manuscript on Lord's gauge theory of gravitation and the brochon: what they mean, why they hold and why they matter},
  pdfauthor={Prepared for Patrick Lee Nash, Ph.D.},
  colorlinks=true,
  linkcolor={blue},
  filecolor={Maroon},
  citecolor={blue},
  urlcolor={blue},
  pdfcreator={LaTeX via pandoc}}

\title{The remarks of the manuscript on Lord's gauge theory of
gravitation and the brochon: what they mean, why they hold and why they
matter}
\usepackage{etoolbox}
\makeatletter
\providecommand{\subtitle}[1]{% add subtitle to \maketitle
  \apptocmd{\@title}{\par {\large #1 \par}}{}{}
}
\makeatother
\subtitle{Provenance document of the remarks, PROJECT spacetimes, folder
Revision\_Lord\_Gravity (task C23, stage 4b)}
\author{Prepared for Patrick Lee Nash, Ph.D.}
\date{November 22, 2026}

\begin{document}
\maketitle

\hypertarget{what-this-document-is}{%
\section{1. What this document is}\label{what-this-document-is}}

The author's standing instruction (here in my words) asks that every
qualifying remark about a result, a remark that says that a result is an
approximation, an estimate, a limit, of a stated accuracy, or restricted
in some way, be explained clearly: what it means, why it holds, and why
it matters. This document keeps these explanations for the manuscript of
the folder, \texttt{lord\_gravity.md}: one section for each remark,
sixteen in all, in the order of the manuscript. Each section says where
the remark is made, states it, and gives the three explanations, the
numbers (written by the programs through the same machine-written
fragments as the manuscript, so that no number is typed twice) and the
checks that decide it. Section 2 explains, once, the objects that the
remarks talk about, so that the document can be read without the
manuscript; Section 19 lists its symbols.

Every statement of this document is a restatement of a statement of the
manuscript or of its programs; it adds no result. ``I'' is the assistant
that did the work.

\hypertarget{the-objects}{%
\section{2. The objects}\label{the-objects}}

\begin{itemize}
\tightlist
\item
  \emph{Units.} \(\hbar=c=1\). The \emph{cutoff} \(\Lambda_{\rm LL}\) is
  the largest momentum of the brochon model; energies are given in units
  of \(\Lambda_{\rm LL}\) unless another unit is named. Newton's
  constant is \(G_{\rm N}=M_{\rm Pl}^{-2}\), with the Planck mass
  \(M_{\rm Pl}\).
\item
  \emph{The brochon} \(\Psi_{\rm LL}\) is a Dirac field with an internal
  index, the isospin, that takes two values (\emph{up} and \emph{down}):
  eight complex components at each point, a Dirac spinor of four
  components for each value of the isospin. Its mass in the normal
  vacuum is \(m_{\rm LL}=0.2\).
\item
  \emph{The brochon model} is a field theory of the brochon with the
  gauge group
  \(SL(2,\mathbb C)\times SU(2)_{\rm v}\times U(1)^{\rm F}_{\rm LL}\)
  (Lorentz transformations of the frame, isospin rotations, phases of
  the fermion number), a contact interaction of four brochon fields, and
  a vacuum in which pairs of brochons of spin zero condense: the
  \emph{superfluid} vacuum. The vacuum without the condensate is the
  \emph{normal vacuum}.
\item
  \emph{The quasi-particle} is the excitation of the brochon in the
  superfluid, a mixture of a brochon and of the complex conjugate of a
  brochon. In the mean-field approximation (each product of field
  operators replaced by its average where the pairs are concerned) it
  has the \emph{Dirac mass} \(M^{\rm CJT,D}_{\rm SF}\) and the
  \emph{pair amplitude} \(D^{\rm CJT}_{\rm SF}\), and the \emph{gap}
  \(E^{\rm CJT,gap}_{\rm SF}=((M^{\rm CJT,D}_{\rm SF})^2+\lvert D^{\rm CJT}_{\rm SF}\rvert^2)^{1/2}=0.3664\),
  the smallest energy of an excitation. The manuscript proves (its
  Theorems L2 and L3) that in every gravitational field and every field
  of the zeugon the quasi-particle behaves exactly as a Dirac particle
  of the mass of the gap.
\item
  \emph{The zeugon} \(A_{{\rm SF},\mu}\) is the gauge boson of
  \(SU(2)_{\rm v}\) that stays massless in the superfluid; it couples to
  the charge \(T^3_{\rm LL}\) (the third isospin generator,
  \(\pm\tfrac12\) for the brochon) with the coupling
  \(g^{\rm v}_{\rm LL}=0.7\). It is not the photon: nothing in the
  brochon model carries electric charge. The other gauge bosons,
  \(W^1\), \(W^2\) and \(B\), are massive in the superfluid.
\item
  \emph{Gravitation.} The gravitational field is described by a
  \emph{tetrad}, four orthonormal vectors at each point (Lord's
  \emph{reference vectors}); the metric, the curvature and the spin
  connection follow from it. The \emph{Schwarzschild} space-time is that
  of a black hole of constant mass \(M_{\rm LG}\) that does not rotate,
  with the \emph{gravitational radius}
  \(\mathcal M_{\rm LG}=G_{\rm N}M_{\rm LG}\) and the horizon
  \(r=2\mathcal M_{\rm LG}\); the \emph{Vaidya} space-time is that of a
  black hole whose mass changes because radiation falls in.
\item
  \emph{The gravitational atom.} Around a black hole a brochon has
  \emph{quasi-bound states}: the levels of a hydrogen atom with the
  coupling \(y_{\rm LG}=G_{\rm N}M_{\rm LG}m_{\rm LL}/(\hbar c)\), which
  decay by falling into the black hole. A state is labelled as in
  hydrogen by its principal number \(n_{\rm LG}\), its orbital angular
  momentum \(\ell_{\rm LG}\) (the letters \(S\), \(P\), \(D\) for 0, 1,
  2) and its total angular momentum \(j_{\rm LG}\), for example
  \(1S_{1/2}\).
\item
  \emph{The pair} is the atom of two quasi-particles of opposite zeugon
  charge, bound by the zeugon with the potential
  \(-\alpha^{\rm Z}_{\rm LG}/r\),
  \(\alpha^{\rm Z}_{\rm LG}=(g^{\rm v}_{\rm LL})^2/(16\pi)=0.00975\): a
  positronium of the mass of the gap.
\end{itemize}

\hypertarget{r1-the-coupling-to-the-scalar-curvature}{%
\section{3. R1: the coupling to the scalar
curvature}\label{r1-the-coupling-to-the-scalar-curvature}}

\emph{Where.} \texttt{lord\_gravity.md}, Section 15.2, after the square
of the Dirac operator in a curved space-time.

\emph{The remark.} The square of the Dirac operator of the brochon
contains the scalar curvature \(R^{\rm LG}\) as a term
\(-\tfrac14R^{\rm LG}\) added to the mass squared, and the field
strengths of the gauge fields through a magnetic moment of gyromagnetic
ratio 2.

\emph{What it means.} A brochon in a curved space-time behaves, to
second order in its derivatives, as a particle whose mass squared is
shifted by a quarter of the scalar curvature:
\((\Box_{\rm LG}+m_{\rm LL}^2-\tfrac14R^{\rm LG}+\dots)\Psi_{\rm LL}=0\),
with the covariant wave operator \(\Box_{\rm LG}\).

\emph{Why it holds.} Squaring the Dirac operator produces the commutator
of two covariant derivatives, which is the curvature of the spin
connection; contracted with the Dirac matrices it gives
\(-\tfrac14R^{\rm LG}\), with the factor \(\tfrac14\) of the
representation of the Lorentz algebra on spinors (the spin connection
acts on a spinor as
\(\tfrac14\omega^{\rm LG}_{\mu\,ab}\gamma^{ab}_{\rm LL}\)). Lord printed
the spinor curvature with \(\tfrac12\) in place of \(\tfrac14\); with
his factor the term would be \(-\tfrac12R^{\rm LG}\), and the checks
exclude it: the deviation from the identity is then as large as the term
itself (relative deviation \(1.000000\), checks G-10b and WG-4).

\emph{Why it matters.} The term acts wherever the scalar curvature is
not zero, inside matter and in the early universe. It does not act in
the space-times of the black holes of the manuscript: the Schwarzschild
space-time has no Ricci curvature, and the Vaidya space-times have
\(R^{\rm LG}=0\) because the energy-momentum tensor of radiation falling
in at the speed of light has no trace (checks WG-1 to WG-3).

\hypertarget{r2-a-constant-dirac-mass-and-pair-amplitude}{%
\section{4. R2: a constant Dirac mass and pair
amplitude}\label{r2-a-constant-dirac-mass-and-pair-amplitude}}

\emph{Where.} Section 15.4 of \texttt{lord\_gravity.md}, after Theorem
L2.

\emph{The remark.} Theorem L2 assumes that the Dirac mass and the pair
amplitude of the quasi-particle are the same everywhere.

\emph{What it means.} The condensate is taken to be that of flat
space-time, unchanged by the gravitational field.

\emph{Why it holds, and when.} The condensate is a property of the
vacuum with the single scale \(\Lambda_{\rm LL}\). A gravitational field
changes it through the curvature, and the relative change is of the
order of the curvature divided by \(\Lambda_{\rm LL}^2\). Near a black
hole of mass \(M_{\rm LG}\) the curvature is at most of the order
\(\mathcal M_{\rm LG}^{-2}\): the invariant
\(R^{\rm LG}_{\mu\nu\rho\sigma}R^{{\rm LG}\,\mu\nu\rho\sigma}\) of the
Schwarzschild space-time is \(48\mathcal M_{\rm LG}^2/r^6\), which at
the horizon is \(3/(4\mathcal M_{\rm LG}^4)\). The corrections are
therefore of the order \((\Lambda_{\rm LL}\mathcal M_{\rm LG})^{-2}\).

\emph{Why it matters.} They are negligible unless the black hole is as
small as the length \(1/\Lambda_{\rm LL}\) of the model. For such a
black hole the condensate itself would have to be computed in the curved
space-time, which this work does not do.

\hypertarget{r3-the-energy-of-the-condensed-vacuum}{%
\section{5. R3: the energy of the condensed
vacuum}\label{r3-the-energy-of-the-condensed-vacuum}}

\emph{Where.} Section 15.4 of \texttt{lord\_gravity.md}; the comparison
with observation in its Section 31.2.

\emph{The remark.} The condensation changes the energy density of the
vacuum, relative to the normal vacuum, by
\(\mathcal E^{\rm CJT}_{\rm SF}\Lambda_{\rm LL}^4\) with
\(\mathcal E^{\rm CJT}_{\rm SF}=-2.422\times10^{-4}\); the brochon model
does not fix the energy of the vacuum to which gravitation responds.

\emph{What it means.} Gravitation responds to a constant energy density
of the vacuum as to a cosmological constant. The model measures every
energy from the normal vacuum (the normal order of its interaction); the
energy that gravitates is this condensation energy plus a constant of
the gravitational theory, which nothing in the model determines.

\emph{Why it holds.} A constant energy density of the vacuum has an
energy-momentum tensor proportional to the metric
\(g^{\rm LG}_{\mu\nu}\), the form of a cosmological constant. Adding a
constant to the Lagrangian changes no field equation of the brochon
model, so the model cannot fix it; only gravitation sees it.

\emph{Why it matters.} The observed energy density of the vacuum is
\(\rho^\Lambda_{\rm LG}=2.52\times10^{-11}\) eV\(^4\) \(=(2.24\)
meV\()^4\) (the Planck satellite, 2018: \(H_0=67.4\) km/s/Mpc,
\(\Omega_{\rm m}=0.315\), a flat universe). The condensation energy has
this size at the cutoff \(\Lambda^*_{\rm LG}=18.0\) meV and exceeds it
by \((\Lambda_{\rm LL}/\Lambda^*_{\rm LG})^4\) above it
(\(9.6\times10^{6}\) for 1 eV, \(9.6\times10^{42}\) for 1 GeV), with the
opposite sign. The constant of the gravitational theory must therefore
cancel the condensation energy to one part in
\((\Lambda_{\rm LL}/\Lambda^*_{\rm LG})^4\), as it must cancel the
energies of the condensates of particle physics: the cosmological
constant problem. (Checks LC-6 and YC-6: the observed density by two
routes, with the Planck mass and with Newton's constant in SI units,
agreeing to \(2.1\times10^{-7}\).)

\hypertarget{r4-operators-normal-order-and-semiclassical-gravitation}{%
\section{6. R4: operators, normal order and semiclassical
gravitation}\label{r4-operators-normal-order-and-semiclassical-gravitation}}

\emph{Where.} Section 15.4 of \texttt{lord\_gravity.md}.

\emph{The remark.} The brochon is a quantum field whose interaction is
normal-ordered with respect to the normal vacuum of flat space-time and
whose momenta are cut off at \(\Lambda_{\rm LL}\); in a curved
space-time neither has a unique generalization. The manuscript states
the field equations of the operators and of the mean field, and
Einstein's equations with the expectation value of the energy-momentum
tensor as the source (\emph{semiclassical gravitation}).

\emph{What it means.} Gravitation is treated classically, the brochon
quantum mechanically, and the gravitational field responds to the
average energy and momentum of the brochon.

\emph{Why it holds.} The split of a field into particles and
antiparticles, which defines the normal vacuum, depends on the time of
an observer when the space-time has no symmetry in time (as the Vaidya
space-times), and a cutoff of momenta must be imposed in a local frame.
Parts IV and V of the manuscript, however, solve one-particle equations
in a given geometry: the brochon is a \emph{test particle}, which needs
neither the normal order nor a renormalized energy-momentum tensor.

\emph{Why it matters.} The results of Parts IV to VI do not depend on
the open questions of quantum fields in curved space-time. The effect of
one brochon on the geometry is smaller than the effects computed by the
ratio of its mass to the mass of the black hole.

\hypertarget{r5-the-massive-gauge-fields}{%
\section{7. R5: the massive gauge
fields}\label{r5-the-massive-gauge-fields}}

\emph{Where.} Section 16.2 of \texttt{lord\_gravity.md}.

\emph{The remark.} In the superfluid \(W^1\), \(W^2\) and \(B\) are
massive; their energy-momentum tensors and that of the pair field add to
the source of Einstein's equations, but the manuscript needs only the
tensor of the zeugon.

\emph{What it means.} Of the four gauge bosons only the zeugon reaches
to infinity. The massive ones act only over a distance of the order of
their inverse masses,
\(1/M^{\rm CJT}_{W,{\rm SF}}\approx23/\Lambda_{\rm LL}\)
(\(M^{\rm CJT}_{W,{\rm SF}}=0.04302319693\)).

\emph{Why it holds.} A massive vector field falls off as \(e^{-Mr}/r\)
around its source; the zeugon, massless, falls off as \(1/r\).

\emph{Why it matters.} Every effect of the gauge fields on the scale of
a black hole, many times \(1/\Lambda_{\rm LL}\), is that of the zeugon:
the charged black hole and the pair of Part VI of the manuscript. The
conservation of the number of quasi-particles (R16) is proved only
without the massive bosons, for which the identity of operators behind
it fails.

\hypertarget{r6-the-formulas-of-the-gravitational-atom-for-small-masses-are-asymptotic}{%
\section{8. R6: the formulas of the gravitational atom for small masses
are
asymptotic}\label{r6-the-formulas-of-the-gravitational-atom-for-small-masses-are-asymptotic}}

\emph{Where.} Section 20.4 of \texttt{lord\_gravity.md}.

\emph{The remark.} For small \(y_{\rm LG}\) the levels and decay rates
follow formulas of S. R. Dolan and D. Dempsey (arXiv:1504.03190, after
Ternov and co-workers): the fine structure
\(F^{\rm H}_{\rm LG}y_{\rm LG}^4\) and the decay rate
\(c^{\rm H}_{\rm LG}y_{\rm LG}^{2\ell_{\rm LG}+2j_{\rm LG}+4}\) (the
power 5 for \(S_{1/2}\), 7 for \(P_{1/2}\), 9 for \(P_{3/2}\), 11 for
\(D_{3/2}\)). They are the leading terms of expansions in
\(y_{\rm LG}\), and at larger \(y_{\rm LG}\) the corrections are large.

\emph{What it means.} The formulas describe the computed states only
when the black hole is small compared with the wavelength of the
brochon, \(y_{\rm LG}\ll1\).

\emph{Why it holds.} The formulas come from matching the hydrogen-like
wave function far from the black hole to the solution near the horizon,
an expansion in the ratio of the two lengths, \(y_{\rm LG}\). The
program computes the states at \(y_{\rm LG}=0.005\) to 0.05 and
extrapolates the ratios of the computed energies and decay rates to the
formulas to \(y_{\rm LG}=0\): they are 1 for all six states, to
\(3.4\times10^{-4}\) and \(7.2\times10^{-4}\) (checks LS-11, LS-12).
Between \(y_{\rm LG}=0.05\) and the next masses the local powers of the
decay rate are \(1S_{1/2}\): 6.03 (\(\mathcal M_{\rm LG}m_{\rm LL}\)
from 0.05 to 0.1); \(2S_{1/2}\): 6.03 (\(\mathcal M_{\rm LG}m_{\rm LL}\)
from 0.05 to 0.1); \(2P_{1/2}\): 7.92 (\(\mathcal M_{\rm LG}m_{\rm LL}\)
from 0.05 to 0.1); \(2P_{3/2}\): 10.87
(\(\mathcal M_{\rm LG}m_{\rm LL}\) from 0.1 to 0.15); \(3P_{3/2}\):
10.90 (\(\mathcal M_{\rm LG}m_{\rm LL}\) from 0.1 to 0.15);
\(3D_{3/2}\): 14.79 (\(\mathcal M_{\rm LG}m_{\rm LL}\) from 0.2 to
0.25), far from the asymptotic ones.

\emph{Why it matters.} The agreement confirms the computed states and
their labels independently. For the states of the gravitational atom
that matter most, near \(y_{\rm LG}\approx0.1\) to 0.4, only the
computed values hold; the formulas only give their trend.

\hypertarget{r7-the-shrinking-black-hole-is-a-model}{%
\section{9. R7: the shrinking black hole is a
model}\label{r7-the-shrinking-black-hole-is-a-model}}

\emph{Where.} Section 23.1 of \texttt{lord\_gravity.md}.

\emph{The remark.} A black hole that shrinks by Hawking's evaporation is
described by the ingoing Vaidya space-time with an inflow of
\emph{negative} energy (the model of W. A. Hiscock, Physical Review D
23, 2813 (1981)); the radiation that leaves the black hole is not
modelled, only its effect at the horizon.

\emph{What it means.} The geometry near the horizon of a shrinking black
hole is represented, the outgoing radiation far away is not.

\emph{Why it holds.} The negative energy that falls in represents the
reaction of the black hole to the emission of Hawking radiation; like
evaporation itself it violates the condition, true for ordinary matter,
that every observer measures a positive energy density. The outgoing
Vaidya space-time, which describes the radiation far away, has
coordinates that are singular at the horizon, so that the field of the
brochon could not be followed across it.

\emph{Why it matters.} The runs with a shrinking mass describe the
brochon near and inside the horizon, where it decays; they do not
describe what a distant observer receives.

\hypertarget{r8-the-projections-on-the-states-of-a-leaking-atom-are-not-probabilities}{%
\section{10. R8: the projections on the states of a leaking atom are not
probabilities}\label{r8-the-projections-on-the-states-of-a-leaking-atom-are-not-probabilities}}

\emph{Where.} Section 27 of \texttt{lord\_gravity.md}.

\emph{The remark.} Around a black hole the quasi-bound states leak and
overlap; the charge of the brochon can be divided among them only
approximately, to about 10 \% when the levels are nearly as broad as
they are apart.

\emph{What it means.} The \emph{occupation} of a state, the squared
cosine of the angle between the field and the state in the charge
product, is not a probability: the occupations of different states need
not add up to at most 1, and the decomposition into states depends on
how it is fitted.

\emph{Why it holds.} The quasi-bound states lose charge into the black
hole, the evolution does not conserve the charge product, and the states
are not orthogonal: the normalized overlap of \(1S_{1/2}\) and
\(2S_{1/2}\) is \(0.135\) at \(y_{\rm LG}=0.2\) and \(0.292\) at 0.25.
At \(y_{\rm LG}=0.25\) the decay rate of \(1S_{1/2}\) is \(0.433\) of
the distance of the levels \(1S_{1/2}\) and \(2S_{1/2}\).

\emph{Why it matters.} The adiabatic and the sudden behaviour of the
runs (Table V.1 of the manuscript) rests on the total charge and on the
fits that agree within 2 \%; the share of a single state is quoted with
this uncertainty.

\hypertarget{r9-the-polarization-of-the-condensate-is-not-computed}{%
\section{11. R9: the polarization of the condensate is not
computed}\label{r9-the-polarization-of-the-condensate-is-not-computed}}

\emph{Where.} Section 29.1 of \texttt{lord\_gravity.md}.

\emph{The remark.} The potential \(-\alpha^{\rm Z}_{\rm LG}/r\) of two
quasi-particles is that of the empty vacuum; the superfluid could weaken
it at long distances by the polarization of its condensate, which is not
computed.

\emph{What it means.} The pair and the charged black hole of Part VI are
computed with the coupling of the vacuum without matter.

\emph{Why it holds.} A medium of charged excitations screens a static
charge. The zeugon stays massless in the superfluid because the
condensate carries no charge \(T^3_{\rm LL}\) and its quasi-particles
have a gap, so that a constant potential of the zeugon changes no charge
(the susceptibility of \(T^3_{\rm LL}\) vanishes); the superfluid
therefore does not screen a static charge completely, as a conductor
would, but it can weaken the field by a dielectric constant different
from 1. Computing it needs the response of the condensate to the zeugon
at small momenta, which this work has not done.

\emph{Why it matters.} A dielectric constant
\(\epsilon^{\rm d}_{\rm LG}\) would replace \(\alpha^{\rm Z}_{\rm LG}\)
by \(\alpha^{\rm Z}_{\rm LG}/\epsilon^{\rm d}_{\rm LG}\) in the pair:
its size would grow by the factor \(\epsilon^{\rm d}_{\rm LG}\) and its
binding shrink by \(1/(\epsilon^{\rm d}_{\rm LG})^2\). The conclusions
of Part VI about the size of the effects (orders of magnitude) would not
change for \(\epsilon^{\rm d}_{\rm LG}\) of order 1.

\hypertarget{r10-the-limit-of-the-charge-by-the-creation-of-pairs-is-an-estimate}{%
\section{12. R10: the limit of the charge by the creation of pairs is an
estimate}\label{r10-the-limit-of-the-charge-by-the-creation-of-pairs-is-an-estimate}}

\emph{Where.} Section 30.3 of \texttt{lord\_gravity.md}.

\emph{The remark.} The charge at which the field at the horizon reaches
the critical field
\((m^{\rm qp}_{\rm LG})^2/\lvert q^{\rm qp}_{\rm LG}\rvert\) of
Schwinger's pair creation (\(m^{\rm qp}_{\rm LG}\) and
\(q^{\rm qp}_{\rm LG}\) the mass and the zeugon charge of a
quasi-particle), \(4y^{\rm qp}_{\rm LG}\) times the limit by accretion
(with
\(y^{\rm qp}_{\rm LG}=G_{\rm N}M_{\rm LG}m^{\rm qp}_{\rm LG}/(\hbar c)\),
\(m^{\rm qp}_{\rm LG}\) the mass of the gives the scale of the limit,
not a sharp value.

\emph{What it means.} Pair creation does not switch on at a threshold;
its rate grows continuously with the field, as
\(\exp(-\pi(m^{\rm qp}_{\rm LG})^2/(\lvert q^{\rm qp}_{\rm LG}\rvert\mathcal F))\)
in a uniform field of strength \(\mathcal F\).

\emph{Why it holds.} The exponent is about \(\pi\) at the critical field
and grows without limit for weaker fields; the discharge becomes
efficient around the critical field, and how efficient depends on the
time available, which the estimate does not fix. (Checks LC-4, YC-4: the
ratio \(4y^{\rm qp}_{\rm LG}\).)

\emph{Why it matters.} The comparison of the two limits, which one is
stronger for \(y^{\rm qp}_{\rm LG}<\tfrac14\), holds within a factor of
order 1; both limits make the charge far too small to be seen.

\hypertarget{r11-the-ringing-is-computed-in-the-limit-of-large-angular-number}{%
\section{13. R11: the ringing is computed in the limit of large angular
number}\label{r11-the-ringing-is-computed-in-the-limit-of-large-angular-number}}

\emph{Where.} Section 30.4 of \texttt{lord\_gravity.md}.

\emph{The remark.} The frequencies of the ringing of a black hole are
related to its shadow by
\(\omega^{\rm QNM}_{\rm LG}\approx(\ell^{\rm QNM}_{\rm LG}+\tfrac12)/b^{\rm sh}_{\rm LG}-i(n^{\rm QNM}_{\rm LG}+\tfrac12)\lambda^{\rm ph}_{\rm LG}\)
only for a large angular number \(\ell^{\rm QNM}_{\rm LG}\) (the
\emph{eikonal} limit; V. Cardoso and co-workers, Physical Review D 79,
064016 (2009)); \(n^{\rm QNM}_{\rm LG}=0,1,\dots\) numbers the modes of
one angular number.

\emph{What it means.} The relative change of the frequencies,
\(+\tfrac16(\mathcal Q^{\rm Z}_{\rm LG}/\mathcal M_{\rm LG})^2\) for a
small zeugon charge \(\mathcal Q^{\rm Z}_{\rm LG}\), is exact for large
\(\ell^{\rm QNM}_{\rm LG}\) and approximate for the modes of small
angular number that dominate a real ringing.

\emph{Why it holds.} For a large angular number a wave circling the
black hole behaves as a ray of light on the photon sphere, whose angular
velocity seen from far away is \(1/b^{\rm sh}_{\rm LG}\); for a small
one the wave is not a ray.

\emph{Why it matters.} For a small angular number the coefficient of
\((\mathcal Q^{\rm Z}_{\rm LG}/\mathcal M_{\rm LG})^2\) in the change of
the frequencies need not be exactly \(\tfrac16\); it was not computed.
At the largest charge the relative change of the shadow is
\(-1.5\times10^{-32}\) for \(\Lambda_{\rm LL}=1\) TeV, so the
conclusion, not detectable, does not depend on the coefficient.

\hypertarget{r12-the-quadratic-tidal-potential-and-the-nonrelativistic-pair}{%
\section{14. R12: the quadratic tidal potential and the nonrelativistic
pair}\label{r12-the-quadratic-tidal-potential-and-the-nonrelativistic-pair}}

\emph{Where.} Section 32.1 of \texttt{lord\_gravity.md}.

\emph{The remark.} The tidal potential of the falling pair is the first
term of the expansion of the tidal field in the size of the pair divided
by the radius, and the relative motion is nonrelativistic.

\emph{What it means.} The pair is treated as a small body: its size, a
few Bohr radii
\(a^{\rm pair}_{\rm LG}=2/(E^{\rm CJT,gap}_{\rm SF}\alpha^{\rm Z}_{\rm LG})\),
is small compared with the radius of the black hole, and its internal
velocities are small compared with the speed of light.

\emph{Why it holds.} At the horizon the ratio of the radius
\(2\mathcal M_{\rm LG}\) to the Bohr radius is
\(2\tilde{\mathcal M}_{\rm LG}/\alpha^{\rm Z}_{\rm LG}\), with the
gravitational radius \(\tilde{\mathcal M}_{\rm LG}\) in the atomic unit
of time of the pair; it exceeds 40 for all the black holes computed
(\(\tilde{\mathcal M}_{\rm LG}\ge0.2\)), so the next term of the
expansion is of relative order \(1/40\) or less. The velocities are of
the order \(\alpha^{\rm Z}_{\rm LG}\approx0.01\).

\emph{Why it matters.} For much lighter black holes the pair would be as
large as the black hole, and the computation would no longer apply; the
conclusion that black holes swallow the pair whole is stated for the
range in which the pair is small.

\hypertarget{r13-the-classical-estimate-of-the-tidal-break-up-is-false-for-a-real-fall}{%
\section{15. R13: the classical estimate of the tidal break-up is false
for a real
fall}\label{r13-the-classical-estimate-of-the-tidal-break-up-is-false-for-a-real-fall}}

\emph{Where.} Sections 32.2 and 32.5 of \texttt{lord\_gravity.md}.

\emph{The remark.} The estimate that black holes lighter than
\(1.492\times10^{5}\) \(M_{\rm Pl}^2/\Lambda_{\rm LL}\) tear the pair
apart before the horizon, from the moment at which the tidal barrier
sinks below the ground level, assumes a static field and fails for a
real fall.

\emph{What it means.} Along the real fall the pair is beyond its barrier
with probability \(\tfrac12\) only for black holes lighter than
\(1.40\times10^{4}\) \(M_{\rm Pl}^2/\Lambda_{\rm LL}\), and even then
its two members cross the horizon within a few Bohr radii of each other
(the probability of a separation beyond 10 Bohr radii is
\(5.2\times10^{-4}\) at the end of the computation).

\emph{Why it holds.} The tidal strength doubles while the remaining time
shrinks by \(2^{-1/2}\), faster than the pair can tunnel (the path under
the barrier takes the imaginary time \(\tfrac34\pi\) times the remaining
time); even without a barrier the escape takes time (rate \(0.0349\) in
the atomic unit when the remaining time is 2); and two particles that
are no longer bound separate only as the remaining time to the power
\(-\tfrac13\). (Checks LT-2 to LT-5, YT-4 to YT-7, WT-5.)

\emph{Why it matters.} The estimate was first recorded as a result. It
is corrected wherever it appears, and the scan over the cutoff no longer
prints it as a prediction.

\hypertarget{r14-the-probability-of-separation-late-in-the-fall-depends-on-a-numerical-device}{%
\section{16. R14: the probability of separation late in the fall depends
on a numerical
device}\label{r14-the-probability-of-separation-late-in-the-fall-depends-on-a-numerical-device}}

\emph{Where.} Section 32.4 of \texttt{lord\_gravity.md} and the report
of the program \texttt{lg\_tidal} (check LT-5b).

\emph{The remark.} Below the remaining time 1 (in the atomic unit of the
pair) the probability that the two quasi-particles are more than 10 Bohr
radii apart depends on the flattening of the tidal potential far out,
which then lies inside the radius 10; at the end of the computation it
lies between \(5.16\times10^{-4}\) and \(7.21\times10^{-4}\) in five
computations with different numerical parameters.

\emph{What it means.} That probability is known there only to some tens
of per cent; the probabilities beyond the barrier and outside the ground
state are not affected (they agree within \(3\times10^{-3}\) among the
computations, check LT-5a).

\emph{Why it holds.} The tidal potential accelerates the escaping waves
without limit; to keep them on the grid the program flattens it to a
constant far beyond the barrier, and the radius of the flattening
shrinks as the field grows. Once it lies inside the radius 10, the
separation beyond 10 measures partly the device.

\emph{Why it matters.} The conclusion, that the members stay within a
few Bohr radii, rests on the root mean squares of the separation
(\(1.78\) Bohr radii along the radius and \(0.77\) across it at the end)
and holds for every one of the five computations.

\hypertarget{r15-the-energy-of-the-vacuum-and-the-values-of-the-cosmological-parameters}{%
\section{17. R15: the energy of the vacuum and the values of the
cosmological
parameters}\label{r15-the-energy-of-the-vacuum-and-the-values-of-the-cosmological-parameters}}

\emph{Where.} Section 31.2 of \texttt{lord\_gravity.md}.

\emph{The remark.} The observed energy density of the vacuum is computed
from the Hubble constant and the matter density of the Planck satellite
in the flat model, with the radiation neglected.

\emph{What it means.} The number
\(\rho^\Lambda_{\rm LG}=2.52\times10^{-11}\) eV\(^4\) carries the
uncertainties of these parameters, about 1.5 \% for \(H_0^2\) and 1 \%
for \(1-\Omega_{\rm m}\) (Planck 2018: \(H_0=67.4\pm0.5\) km/s/Mpc,
\(\Omega_{\rm m}=0.315\pm0.007\)); the radiation, about \(10^{-4}\) of
the critical density, changes \(\Omega_\Lambda\) by less.

\emph{Why it holds.}
\(\rho^\Lambda_{\rm LG}=(1-\Omega_{\rm m})\,3H_0^2/(8\pi G_{\rm N})\) in
the flat model.

\emph{Why it matters.} An uncertainty of a few per cent in
\(\rho^\Lambda_{\rm LG}\) changes \(\Lambda^*_{\rm LG}\) by its fourth
root, less than 1 \%. The comparison of the condensation energy with the
observed one is a comparison of orders of magnitude, and it is
unaffected.

\hypertarget{r16-the-conservation-of-the-number-of-quasi-particles-holds-in-the-mean-field}{%
\section{18. R16: the conservation of the number of quasi-particles
holds in the mean
field}\label{r16-the-conservation-of-the-number-of-quasi-particles-holds-in-the-mean-field}}

\emph{Where.} Section 33.4 of \texttt{lord\_gravity.md}.

\emph{The remark.} The proof that the zeugon conserves the number of
quasi-particles \(N^{\rm qp}_{\rm LG}\) (the particles minus the
antiparticles of the Dirac field of the mass of the gap), and with it
the rule that two quasi-particles of opposite charge cannot annihilate
into zeugons, holds for the mean-field equation of the superfluid and
for the zeugon alone.

\emph{What it means.} Two corrections are left out: the part of the
four-fermion interaction that the mean field does not keep, which
conserves the number of brochons but not \(N^{\rm qp}_{\rm LG}\) and can
change it by 2 at a higher order; and the massive gauge bosons \(W^1\),
\(W^2\) and \(B\), for which the identity of operators behind the proof
fails.

\emph{Why it holds.} The quasi-particle field is
\(\Psi_{\rm LL}=\chi_{\rm LG}+b^{\rm qp}_{\rm LG}\mathsf P_{\rm LG}\chi_{\rm LG}^{\dagger T}\)
with a Dirac field \(\chi_{\rm LG}\) of the mass of the gap (a
consequence of the mean-field equation); its zeugon current is
\((1+\lvert b^{\rm qp}_{\rm LG}\rvert^2)\) times that of
\(\chi_{\rm LG}\), plus a constant, because the terms that would change
\(N^{\rm qp}_{\rm LG}\) are symmetric combinations of anticommuting
operators (checks AN-3a, AN-3b, YA-3, WA-3, to \(4.5\times10^{-16}\)).
The interaction beyond the mean field contains products of four brochon
fields that need not have this property.

\emph{Why it matters.} The pairs of two particles
(\(N^{\rm qp}_{\rm LG}=\pm2\)) are stable against annihilation into
zeugons only to the order computed; their annihilation through the
residual interaction is suppressed by a higher power of the couplings,
at a rate not computed.

\hypertarget{the-symbols-of-this-document}{%
\section{19. The symbols of this
document}\label{the-symbols-of-this-document}}

The symbols are those of the manuscript \texttt{lord\_gravity.md}; this
table lists those that this document uses.

\begin{longtable}[]{@{}
  >{\raggedright\arraybackslash}p{(\columnwidth - 2\tabcolsep) * \real{0.5000}}
  >{\raggedright\arraybackslash}p{(\columnwidth - 2\tabcolsep) * \real{0.5000}}@{}}
\toprule\noalign{}
\begin{minipage}[b]{\linewidth}\raggedright
symbol
\end{minipage} & \begin{minipage}[b]{\linewidth}\raggedright
meaning
\end{minipage} \\
\midrule\noalign{}
\endhead
\bottomrule\noalign{}
\endlastfoot
\(\hbar\), \(c\), \(G_{\rm N}\), \(M_{\rm Pl}\) & Planck's constant, the
speed of light, Newton's constant and the Planck mass
\((\hbar c/G_{\rm N})^{1/2}\) \\
\(\Lambda_{\rm LL}\) & the momentum cutoff of the brochon model, the
unit of energy \\
\(\Psi_{\rm LL}\), \(m_{\rm LL}\) & the brochon and its mass in the
normal vacuum \\
\(M^{\rm CJT,D}_{\rm SF}\), \(D^{\rm CJT}_{\rm SF}\),
\(E^{\rm CJT,gap}_{\rm SF}\) & the Dirac mass, the pair amplitude and
the gap of the superfluid (method CJT) \\
\(A_{{\rm SF},\mu}\), \(W^1\), \(W^2\), \(B\) & the zeugon and the
massive gauge bosons \\
\(g^{\rm v}_{\rm LL}\), \(T^3_{\rm LL}\) & the coupling of the zeugon
and the charge to which it couples \\
\(M_{\rm LG}\), \(\mathcal M_{\rm LG}\), \(r\) & the mass of a black
hole, its gravitational radius \(G_{\rm N}M_{\rm LG}\), and the radius
(or, for the pair, the distance of its members) \\
\(y_{\rm LG}\), \(y^{\rm qp}_{\rm LG}\) &
\(G_{\rm N}M_{\rm LG}m_{\rm LL}/(\hbar c)\) and
\(G_{\rm N}M_{\rm LG}m^{\rm qp}_{\rm LG}/(\hbar c)\), the couplings of
the gravitational atom of the brochon and of the quasi-particle \\
\(m^{\rm qp}_{\rm LG}\), \(q^{\rm qp}_{\rm LG}\), \(\mathcal F\) & the
mass and the zeugon charge of a quasi-particle, and the strength of a
uniform field in Schwinger's rate \\
\(n_{\rm LG}\), \(\ell_{\rm LG}\), \(j_{\rm LG}\) & the principal
number, the orbital and the total angular momentum of a state of the
gravitational atom \\
\(F^{\rm H}_{\rm LG}\), \(c^{\rm H}_{\rm LG}\) & the coefficients of the
fine structure and of the decay rate of the gravitational atom at small
\(y_{\rm LG}\) \\
\(R^{\rm LG}\), \(R^{\rm LG}_{\mu\nu\rho\sigma}\),
\(g^{\rm LG}_{\mu\nu}\), \(\Box_{\rm LG}\),
\(\omega^{\rm LG}_{\mu\,ab}\), \(\gamma^{ab}_{\rm LL}\) & the scalar
curvature, the Riemann tensor, the metric, the covariant wave operator,
the spin connection and the antisymmetrized Dirac matrices \\
\(\alpha^{\rm Z}_{\rm LG}\), \(a^{\rm pair}_{\rm LG}\) & the coupling of
the zeugon between two quasi-particles and the Bohr radius of the
pair \\
\(\tilde{\mathcal M}_{\rm LG}\) & the gravitational radius of a black
hole in the atomic unit of time of the pair \\
\(\mathcal E^{\rm CJT}_{\rm SF}\) & the condensation energy per volume
in units of \(\Lambda_{\rm LL}^4\) \\
\(\rho^\Lambda_{\rm LG}\), \(\Lambda^*_{\rm LG}\), \(H_0\),
\(\Omega_{\rm m}\), \(\Omega_\Lambda\) & the observed energy density of
the vacuum, the cutoff at which the condensation energy has its size,
the Hubble constant and the density parameters \\
\(\mathcal Q^{\rm Z}_{\rm LG}\), \(b^{\rm sh}_{\rm LG}\),
\(\lambda^{\rm ph}_{\rm LG}\) & the zeugon charge of a black hole as a
length, the radius of its shadow, and the rate at which rays leave the
photon sphere \\
\(\omega^{\rm QNM}_{\rm LG}\), \(\ell^{\rm QNM}_{\rm LG}\),
\(n^{\rm QNM}_{\rm LG}\) & the complex frequency, the angular number and
the number of a mode of the ringing \\
\(\epsilon^{\rm d}_{\rm LG}\) & a dielectric constant of the condensate
(Section 11 of this document) \\
\(N^{\rm qp}_{\rm LG}\), \(\chi_{\rm LG}\), \(b^{\rm qp}_{\rm LG}\),
\(\mathsf P_{\rm LG}\) & the number of quasi-particles, the Dirac field
of the mass of the gap, the coefficient \(-D/(E+M)\), and the matrix of
the antilinear map of the pairing \\
\end{longtable}

\end{document}
